% =============================================================
% Section: Introduction
% =============================================================

\section{Introduction}\label{sec:intro}

Brazil is debating an eight-hour cut to its standard workweek. PEC 8/2025 would lower the cap from 44 to 36 hours and bind roughly two thirds of formal employees at once. Two numbers anchor the policy discussion. The first is the output removed on impact. The second is the one-off productivity gain required to restore it.

The Brazilian setting complicates both answers. Informality is $37.8\%$ under the narrow PNAD definition and $44.2\%$ under the broad IBGE one, ranging from $41\%$ in services to $72\%$ in agriculture. Standard treatments of hours reform model a fully formal economy \citep{Cacciatore2016} or treat informality as a static size distortion \citep{Ulyssea2018}. Neither has been applied to a large cap cut.

This paper builds a static partial-equilibrium model disciplined by Brazilian microdata. A CES aggregator over formal and informal effective labor governs reallocation. A fatigue function $e(h)$ with peak at $40$ hours maps the cap into per-hour productivity. Two firm-size groups generate differential incidence and GHH preferences support welfare. The substitution elasticity is pinned to the Brazilian formal--informal wage premium of \citet{MeghirNaritaRobin2015}, giving $\sigma_{\text{sub}}=1.326$ with interval $[1.116,\,1.469]$ at $(\omega,\eta_I)=(0.622,\,0.40)$.

The central object is $A_{\text{req}}$, the one-off TFP gain that restores pre-reform output. Under the preferred flat-below specification $A_{\text{req}}=6.63\%$. Under the conservative symmetric specification $A_{\text{req}}=8.18\%$. The joint $(\sigma,\omega,\eta_I)$ envelope is $[5.62\%,\,7.48\%]$ under the preferred specification and $[6.93\%,\,9.23\%]$ under the conservative one. At Brazil's best post-1990 decade (PWT \texttt{rtfpna}, $+0.26\%$ per year, 2003--2013) accumulating $A_{\text{req}}$ takes $25$--$31$ years. A stress test that layers a $1$--$2\%$ post-reform TFP contraction raises the conservative $A_{\text{req}}$ to $9.28$--$10.39\%$.

Aggregate informality rises by $+1.6$--$1.9$ pp. With $\sigma>1$ and $\eta_I=0.40$, formal and informal labor are imperfect substitutes, so the cap pushes marginal formal workers toward the informal sector instead of being absorbed by intensive adjustment alone. Welfare peaks near a $41$-hour cap and turns negative below $38$.

The paper makes three contributions. First, it embeds a formal--informal reallocation margin into a structural analysis of hours reform and shows the margin is modest under wage-premium-disciplined parameters. Second, it proposes $A_{\text{req}}$ as a benchmark commensurable with the historical TFP record. Third, it provides a wage-premium-conditional interval for $\sigma_{\text{sub}}$ that the reader can reassess.

Scope is limited. The framework is static and partial equilibrium. It omits capital deepening, general-equilibrium price adjustment, worker selection, and work intensification. Evidence from Portugal 1996 \citep{AsaiLopesTondini2024} and Brazil 1988 \citep{GonzagaMenezesFilhoCamargo2003} enters as a directional plausibility check rather than an identified lower bound for Brazil's effective $A_{\text{req}}$.

Section~\ref{sec:framework} presents facts, model, and calibration. Section~\ref{sec:validation} validates the fit. Section~\ref{sec:results} reports the counterfactual, the TFP benchmark, and the stress test. Section~\ref{sec:conclusion} draws policy implications.
